% !TeX root = ../main.tex

\chapter{基于问题条件区域—属性图增强的交通场景视觉问答方法}

\section{引言}

\section{交通场景视觉问答问题定义与总体框架}
\subsection{问题、视觉观测与外部知识表示}
\subsection{交通场景关系理解与知识推理任务}
\subsection{PerGraph 原始方法与问答扩展的对应关系}

\section{问题条件的动态区域图构建}
\subsection{视觉区域与问题语义编码}
\subsection{问题相关区域节点激活}
\subsection{空间与语义关系构建}

\section{问题条件的语义属性图构建}
\subsection{交通场景属性与知识表示}
\subsection{属性关系构建与问题相关性建模}
\subsection{语义属性节点动态激活}

\section{区域—属性跨图推理与答案预测}
\subsection{问题引导的双向跨图交互}
\subsection{区域—属性一致性剪枝}
\subsection{关系信息聚合与答案预测}
\subsection{训练目标与优化策略}

\section{实验结果与分析}
\subsection{交通场景视觉问答数据与标注组织}
\subsection{关系推理问答样本构建与数据划分}
\subsection{原始感知评价与交通问答的验证设置}
\subsection{对比方法与评价指标}
\subsection{交通场景视觉问答性能分析}
\subsection{区域图、属性图与问题条件消融}
\subsection{跨图交互与一致性剪枝消融}
\subsection{关系推理与跨场景泛化分析}
\subsection{典型问答与失败案例分析}

\section{本章小结}
